% FILE: 03_architecture_and_primitives.tex
% Project: research-prompt-manufactory
% Thesis Role: prompt-manufacturing-and-subagent-command-surface

\section{Architecture and Primitives}
\label{sec:research-prompt-manufactory-architecture}

\subsection{Core Objects}
\label{sec:research-prompt-manufactory-architecture-objects}

The Prompt Manufactory's object model is defined in \texttt{prompt-manufactory.ttl} and
comprises 15 classes grounded in public vocabularies (PROV-O, DCTERMS, DCAT, SKOS, SHACL).
The core objects and their roles in the manufacturing pipeline are:

\begin{description}
  \item[\texttt{pm:ResearchProgram}] The top-level program node. Subclass of
    \texttt{prov:Entity}. Carries a \texttt{pm:programId} (required by SHACL), a
    \texttt{pm:mission} string, and \texttt{pm:hasWorkflow} references. Seven seed instances
    are encoded in \texttt{research-program-law.ttl}: PI, GGEN\_ECOSYSTEM, GGEN\_OTEL,
    ZOEAPP, EXPO\_SUPABASE, CLAUDE\_WORKFLOW, WASM4PM\_REMEDIATE.

  \item[\texttt{pm:Workflow}] A phase-sequenced execution container. The INTEL workflow
    defines 8 ordered phases: Census, Classify, Manifest, Queries, Templates, Conformance,
    Reconciliation, and Checkpoint. The REMEDIATE workflow defines 1 phase. Both are encoded
    in \texttt{workflow-law.ttl} as named RDF instances with \texttt{pm:hasPhase} links.

  \item[\texttt{pm:Phase}] A bounded stage within a workflow. Carries a label and mission,
    and links to one or more \texttt{pm:SubagentRole} instances via
    \texttt{pm:hasSubagentRole}. Phase 1 (Census) assigns 7 subagent roles simultaneously;
    subsequent phases assign 1 each.

  \item[\texttt{pm:SubagentRole}] A bounded inspection station. Required by SHACL to carry
    \texttt{pm:mission}. Also carries \texttt{pm:ownsSurface} (file or directory the
    subagent is authorized to inspect), \texttt{pm:forbidsSurface} (surface the subagent
    must not touch), \texttt{pm:hasOutputContract} (required output artifact path), and
    \texttt{pm:hasRefusalGate} (condition under which the subagent must refuse to proceed).
    Fifteen roles are defined in \texttt{subagent-role-law.ttl}.

  \item[\texttt{pm:HookPolicy}] An Andon gate --- a deterministic lifecycle enforcement
    hook. Six policies are defined in \texttt{hook-law.ttl}, governing conditions under
    which the manufacturing pipeline must halt before proceeding.

  \item[\texttt{pm:Checkpoint}] An ALIVE/PARTIAL verdict record. Required by SHACL to
    carry \texttt{dct:issued} with a \texttt{xsd:dateTime} value. Ten gate definitions
    are encoded in \texttt{checkpoint-law.ttl}.

  \item[\texttt{pm:RenderedPrompt}] The manufactured research warrant. Required by SHACL
    to carry \texttt{pm:derivedFrom} pointing to at least one \texttt{prov:Entity} --- the
    ontology source from which the warrant was manufactured. This constraint is the
    foundational proof-of-manufacture requirement.

  \item[\texttt{pm:Receipt}] A proof-of-execution artifact. Carries a
    \texttt{pm:proves} link to the \texttt{pm:RenderedPrompt} it certifies and a derivation
    hash. The receipt is the replay anchor.

  \item[\texttt{pm:PromptClass}] A type enumeration for warrants:
    INTEL $\mid$ RECONCILE $\mid$ REMEDIATE $\mid$ SYNC $\mid$ AUDIT $\mid$ CHECKPOINT
    $\mid$ SPR $\mid$ LIVESTREAM. Research programs declare their prompt class via
    \texttt{pm:hasPromptClass}.

  \item[\texttt{pm:InvalidGgenFile}] A classification record for legacy \texttt{.ggen}
    files. Carries \texttt{pm:filePath}, \texttt{pm:ownerProject},
    \texttt{pm:classification}, \texttt{pm:blockingStatus}, and
    \texttt{pm:remediationRoute}. Twenty-two instances are encoded in
    \texttt{forbidden-collapse-law.ttl}.
\end{description}

\subsection{Admissible Transitions}
\label{sec:research-prompt-manufactory-architecture-transitions}

The Prompt Manufactory enforces a strict law-to-artifact transition pipeline. Admissible
transitions in the manufacturing graph are:

\begin{enumerate}
  \item \textbf{Graph loading}: All 8 \texttt{.ttl} ontology files load into the RDF graph
    without parse error. This is the entry gate.
  \item \textbf{SPARQL extraction}: A \texttt{.rq} query executes against the loaded graph
    and returns a non-empty result set binding the variables required by its paired template.
  \item \textbf{Tera rendering}: The result set bindings are injected into the \texttt{.tera}
    template, producing a Markdown artifact at the path specified by the generation rule's
    \texttt{output\_file} field.
  \item \textbf{Derivation assertion}: The rendered artifact receives a
    \texttt{pm:derivedFrom} triple in the graph, asserting its provenance.
  \item \textbf{Receipt emission}: An entry is created in the \texttt{prompt-receipt-ledger.md}
    recording the artifact's provenance chain, manufacture timestamp, and proof gates.
  \item \textbf{Manifest recording}: The manufacturing manifest
    (\texttt{manufacturing-manifest-20260601.yaml}) records the total artifacts emitted and
    the gate check results.
\end{enumerate}

Transitions are governed by ggen v26.5.21, which executes 6 generation rules in the order
specified in \texttt{ggen.toml}. Conflict mode is set to \texttt{fail}, meaning any
generation rule that produces a conflict with an existing artifact halts the pipeline rather
than silently overwriting.

\subsection{Boundaries}
\label{sec:research-prompt-manufactory-architecture-boundaries}

The Prompt Manufactory enforces strict surface boundaries. The source layer (\texttt{ggen/})
accepts only \texttt{.ttl}, \texttt{.rq}, \texttt{.tera}, and \texttt{ggen.toml} files.
The emission layer (\texttt{emitted/}) is the only valid destination for manufactured
artifacts. Subagent roles are bounded by \texttt{pm:ownsSurface} and
\texttt{pm:forbidsSurface} assertions: the Census role for the ggen surface
(\texttt{role\_ggen\_census}) explicitly forbids \texttt{.ggen} files as source surfaces;
the engine census role forbids \texttt{../wasm4pm-compat/} and
\texttt{../process-intelligence/}.

The \texttt{ggen.toml} records \texttt{hand\_coding = false} as a project-level assertion,
establishing that no artifact in the emission layer may have been authored by hand. Any
artifact lacking a \texttt{pm:derivedFrom} triple violates the SHACL shape and fails the
``no hand-written program prompts'' audit gate.

\subsection{Failure Modes Refused}
\label{sec:research-prompt-manufactory-architecture-failures}

The \texttt{forbidden-collapse-law.ttl} codifies three failure modes that the Prompt
Manufactory absolutely refuses:

\begin{description}
  \item[\texttt{pm:FORBIDDEN\_ggen\_source\_files}] No \texttt{.ggen} extension files
    permitted in source directories. Audit rule: \texttt{find . -name `*.ggen'} in
    \texttt{prompt-manufactory/ggen/} must return zero results. As of checkpoint
    PARTIAL\_001, this gate passes.

  \item[\texttt{pm:FORBIDDEN\_dto\_flattening}] No collapse of structured
    \texttt{Evidence<T, State, W>} into casual JSON/string carriers without boundary
    classification. Examples include \texttt{EvidenceDto}, \texttt{payload\_json}, or
    \texttt{state\_tag} as unbounded strings. Refusal condition:
    \texttt{REFUSAL\_DTO\_COLLAPSE}.

  \item[\texttt{pm:FORBIDDEN\_forced\_alive}] No ALIVE verdict claimed when critical
    audit gates have failed or prerequisite proofs are incomplete. The checkpoint record
    explicitly states: \textit{``Declaring ALIVE from file-count completion alone would be
    propaganda. PARTIAL acknowledges what works and what awaits.''}
\end{description}

These refusals are encoded as named RDF classes, not as documentation commentary, making
them queryable, auditable, and enforceable through SPARQL ASK queries that the manufacturing
pipeline can execute as proof gates before emitting any ALIVE verdict.
